\scanpage{69}
% loc: 1–6, top
\uline{Remarque} Indépendamment de tte hypo sur $\operatorname{Pic}(k)$ et sur $M$, on peut regarder $C^*$, et la structure de graphe sur $C^\times$ définie par la relation $\operatorname{Tr}(uv)=-1$, d'où structure triangulée. Mais

% loc: 7–9
1\textsuperscript{o}) $C^*$ \deleted{\unsure{ne|n'est}} n'est non vide que si $\deleted{\unsure{k}}M$ est triviale, plus précisément si $M\simeq\operatorname{End}(E)$ avec $E=k\oplus L$ ($L$ module inv.)

% loc: 10–17
\deleted{On a dit} 2\textsuperscript{o}) \deleted{Il n'a \unsure{Il|Je}} L'ens des arêtes est non vide ssi $M\simeq M_2(k)$. En effet, si $\exists$ arête non vide $\{u,v\}$ \deleted{$\operatorname{Tr}uv=$} avec $u,v\in C^\times$, $\operatorname{Tr}uv=-1$, les théorèmes de conjugaison nous fournissent un iso unique $\mathbb M_2\longrightarrow M$, transformant $u_0$ en $u$, $v_0$ en $v$ ($(u_0,v_0)$ couple type dans $M$, satisfaisant les $\unsure{={}^s|\text{mêmes}|=}$ conditions \unsure{:|.}).

% loc: 18–20, bottom
Donc on a intérêt à se borner \unsure{alors|tjs} à des iso $M\simeq\operatorname{End}(E)$, en supposant $E$ \uline{libre}
% archive sheet number: 68

